\noindent\textbf{Background:} Human genetics can prioritize therapeutic mechanisms, but clinically validated pathways may be untestable when strong tissue-relevant instruments are absent. We quantified cis-eQTL coverage and screened 19 dermatologic immune pathway genes against four inflammatory skin-disease outcomes.

\noindent\textbf{Methods:} Using GTEx v8 significant-pair files, we selected one lead cis-eQTL per gene--tissue cell in sun-exposed skin, non-sun-exposed skin, and whole blood. Primary instruments required \textit{p} $<5\times10^{-8}$ and F $>10$; relaxed instruments (\textit{p} $<10^{-5}$) were analyzed separately. Single-variant Wald ratios for FinnGen R13 psoriasis, dermatitis and eczema, vitiligo, and alopecia areata were FDR-corrected within each testing family. TYK2 signals underwent external association lookup and Bayesian colocalization.

\noindent\textbf{Results:} Primary instruments were available for 18 of 57 cells (31.6\%); one PDE4A whole-blood variant was absent from all outcomes, leaving 68 primary tests. Three FDR-significant results linked higher genetically proxied TYK2 expression to lower psoriasis risk in sun-exposed skin (OR 0.454, 95\% CI 0.371--0.557), non-sun-exposed skin (OR 0.604, 95\% CI 0.510--0.716), and whole blood (OR 0.507, 95\% CI 0.403--0.637). These three tissue estimates used two variants because rs280497 was shared by non-sun-exposed skin and whole blood. All 23 exact-variant estimates across eight external GWAS records were protective. Colocalization supported a shared TYK2--psoriasis signal in sun-exposed skin (PP.H4 0.825) and whole blood (PP.H4 0.782), but not non-sun-exposed skin (PP.H4 0.220). Relaxed-threshold coverage reached 23 of 57 cells; the PDE4A--psoriasis locus-level analysis was non-supportive (PP.H4 0.108), but the MR SNP was absent from the colocalization rows.

\noindent\textbf{Conclusion:} The screen recovered the established TYK2 psoriasis locus while showing that instrument scarcity constrains dermatologic drug-target MR. Transcript direction should not be equated with pharmacologic inhibition, absence of an instrument is not evidence against a target, and the PDE4A result remains exploratory.

\noindent\textbf{Keywords:} Mendelian randomization, cis-eQTL, psoriasis, TYK2, colocalization, GTEx